\documentclass[11pt]{article}
\usepackage[a4paper,margin=18mm]{geometry}
\usepackage{amsmath,amssymb}
\usepackage[T1]{fontenc}
\usepackage{lmodern}
\usepackage{microtype}
\usepackage{enumitem}
\setlength{\parindent}{0pt}
\setlength{\parskip}{4pt}
\setlist[itemize]{leftmargin=5mm,itemsep=1.7pt,topsep=3pt}
\pagestyle{empty}

\begin{document}

{\Large\bfseries Trajectory Preprocessing}\par

Each sample contains 30 observed vessel positions and 30 future positions,
\[
H=(x_1,\ldots,x_{30}), \qquad
Y=(y_1,\ldots,y_{30}),
\]
where every \(x_t,y_t\in\mathbb{R}^2\).

The transformed representation describes the trajectory relative to the vessel's \textbf{last observed position} and \textbf{latest observed direction of motion}. Both references are obtained from the observed history \(H\) and are then used unchanged for the future \(Y\).

The last observed position is \(x_{30}\). Every position \(p\) is expressed relative to it as
\[
p-x_{30}.
\]
The observation therefore ends at
\[
x_{30}-x_{30}=(0,0),
\]
and future positions describe where the vessel moves relative to the point where observation ended.

Compass orientation is removed in the same way. The vessel's latest observed direction is obtained from consecutive history positions,
\[
\Delta x_t=x_t-x_{t-1},
\qquad t=2,\ldots,30.
\]
Starting from the end of the history and moving backward, we take the first displacement satisfying
\[
\|\Delta x_t\|_2>10^{-8}.
\]
Its heading is
\[
\theta=
\operatorname{atan2}
(\Delta x_{t,y},\Delta x_{t,x}).
\]
If no such displacement exists, \(\theta=0\).

The trajectory is rotated so this direction points along the positive \(x\)-axis:
\[
\tilde p=R(-\theta)(p-x_{30}),
\]
where
\[
R(\phi)=
\begin{bmatrix}
\cos\phi & -\sin\phi\\
\sin\phi & \cos\phi
\end{bmatrix}.
\]
After this transformation, every trajectory is viewed from the same position and direction: the observation ends at \((0,0)\), with the vessel facing along the positive \(x\)-axis.

The aligned positions can be used directly. Alternatively, they can be represented by how the vessel moves from one aligned point to the next. For an aligned trajectory
\[
(q_1,\ldots,q_{30}),
\]
we store
\[
d_1=q_1,
\qquad
d_t=q_t-q_{t-1},
\quad t=2,\ldots,30.
\]
The first point is kept so no positional information is lost. For the future trajectory, \(d_1\) is the first future position relative to \(x_{30}\); every following value is the movement between two consecutive future positions.

This gives the three representations used by the experiment system: \textbf{raw coordinates}, \textbf{heading-aligned positions}, and \textbf{heading-aligned displacements}.

The transformation is reversible. Displacements recover aligned positions through
\[
q_1=d_1,
\qquad
q_t=q_{t-1}+d_t,
\]
and aligned positions return to the original coordinates through
\[
p=R(\theta)q+x_{30}.
\]

Normalization is applied after this representation is formed and is handled separately, since ordinary trajectory models and chart-based flow models use different normalization layouts.

\newpage
{\large\bfseries Specs}\par

\begin{itemize}
\item Observed trajectory: \textbf{30 positions}
\item Future trajectory: \textbf{30 positions}
\item Position dimension: \textbf{2}
\item Sampling interval: \textbf{20 s}
\item Last observed reference point: \(\mathbf{x_{30}}\)
\item Consecutive history displacements available for heading: \textbf{29}
\item Minimum displacement norm accepted as motion: \(\mathbf{10^{-8}}\)
\item Heading fallback when all observed displacements are below the threshold: \(\mathbf{0}\) rad
\item Forward alignment rotation: \(\mathbf{-\theta}\)
\item Inverse rotation: \(\mathbf{+\theta}\)
\item Aligned final observed position: \(\mathbf{(0,0)}\)
\item Heading-aligned displacement sequence: \textbf{30 two-dimensional values}
\item Displacement sequence contents: \textbf{1 initial aligned position + 29 consecutive displacements}
\item Representation shape before any chart conversion: \(\mathbf{30\times2}\)
\item Raw representation shape: \(\mathbf{30\times2}\)
\item Preprocessing output dtype in the implementation: \textbf{float32}
\end{itemize}

\end{document}
